\section{Jet cones and the canonical interior point}
\label{jet:section}

We work with the fixed cone supplied by
Proposition~\ref{cone:reduction}. Thus $C=\Spec S$ has dimension
$n=N+1$, its puncture is smooth, and its canonical module has a
homogeneous nowhere-vanishing generator on the puncture. The minimizing
Reeb vector $\xi$ is normalized by $A_C(\xi)=n$. Throughout this section,
\begin{equation}\label{jet:density-hypothesis}
 \nvol(o,C)\ge 2N^n.
\end{equation}
Take the primitive positive gradings of
Lemma~\ref{cone:approximation}, so that
\[
 \zeta_\nu=(r_\nu/n)\xi+\epsilon_\nu,
 \qquad r_\nu=A_C(\zeta_\nu),\qquad \epsilon_\nu\longrightarrow0
\]
in absolute cocharacter norm. In the rational case we use the exact
constant grading. Consider a subsequence on which the largest stabilizer
order satisfies $m_\nu\ge r_\nu/N$, and choose a point attaining that
order. We suppress the index $\nu$ when discussing a single grading.
Constants carrying a subscript $C$ may depend on this fixed cone and on
$N$; no uniformity over different cones is needed.

\subsection{Leading Taylor terms and the real cone}

Let $B$ and $L$ be the smooth quotient orbifold and its weight-one line
bundle from Lemma~\ref{cone:orbifold}. In particular,
\[
 S_j=H^0(B,jL),\qquad -K_B=rL.
\]
At the chosen point, take a smooth equivariant uniformizer with parameters
$z_1,\ldots,z_N$ and stabilizer $\mu_m$. Choose the parameters and the
frame of $L$ so that their characters are respectively
$b_1,\ldots,b_N\in\ZZ/m$ and $1$. We choose these algebraically:
averaging regular lifts of a cotangent eigenbasis gives an equivariant
map to $\A^N$ that is étale at the chosen point, and averaging a
nonvanishing local frame gives the specified line character.
After shrinking, these are equivariant algebraic étale coordinates
and an algebraic frame. A section of $jL$ is then represented
by a power series $F(z)$ satisfying
\[
 F(\varepsilon^{b_1}z_1,\ldots,\varepsilon^{b_N}z_N)
 =\varepsilon^j F(z).
\]
Order monomials first by total degree and then lexicographically, with
the smallest monomial first. Write $\operatorname{in}(s)=z^\gamma$ for
the leading monomial, omitting its nonzero scalar coefficient, and put
\begin{align}
 \Gamma&=\{(j,\gamma):0\ne s\in S_j,
                     \ \operatorname{in}(s)=z^\gamma\},
                     \label{jet:semigroup}\\
 \Lambda&=\{(j,\gamma)\in\ZZ\times\ZZ^N:
                  j\equiv\textstyle\sum_i b_i\gamma_i\pmod m\},
                     \label{jet:lattice}\\
 K&=\overline{\operatorname{cone}(\Gamma)}
                       \subset\RR\times\RR^N.
                     \label{jet:cone}
\end{align}
Here $(0,0)\in\Gamma$, and $\gamma=0$ denotes a nonzero constant Taylor
term, including those of sections of positive degree. The degree
coordinate in Equations~\eqref{jet:semigroup}--\eqref{jet:cone} is the
\emph{original} degree $j$.

For integrations only, introduce the scaled degree and the number
\begin{equation}\label{jet:scaling}
 s=\frac r{Nm}\le1,\qquad
 D_0=\frac{j}{sm}=\frac{Nj}{r}.
\end{equation}
We distinguish the coordinate $D_0$ from $j$ throughout. Write
$\widetilde K$, $\widetilde\Gamma$, and $\widetilde\Lambda$ for the
images of $K$, $\Gamma$, and $\Lambda$ under
$(j,x)\mapsto(Nj/r,x)$.

\begin{proposition}[Jet setup]\label{jet:setup}
The following assertions hold.
\begin{enumerate}[label=\textup{(\roman*)}]
\item
$\Gamma$ is an additive semigroup with one-dimensional leaves, and its
generated group is exactly $\Lambda$. Every finite compatible Taylor
jet is realized by sections of all sufficiently large degrees in each
admissible progression.
\item
$K$ is a full-dimensional closed convex cone, contains the positive
constant ray, projects onto the whole exponent orthant, and satisfies
$K=\overline{\Int K}$. Its bounded-degree sections are bounded. More
precisely, uniformly in the chosen point and nearby grading,
\begin{equation}\label{jet:bezout-bound}
 |\gamma|:=\sum_i\gamma_i\le C_C\frac jr
                       \qquad ((j,\gamma)\in K).
\end{equation}
Its scaled image has the form
\begin{equation}\label{jet:epigraph}
 \widetilde K=\{(D_0,x):x\in\RR_{\ge0}^N,\ D_0\ge f(x)\},
\end{equation}
where $f$ is finite, positively
homogeneous, convex, and
\begin{equation}\label{jet:f-coercive}
 f(x)\ge c_C|x|\qquad(x\in\RR_{\ge0}^N)
\end{equation}
for a constant $c_C>0$ independent of the nearby grading and chosen
point.
\item
In the scaled coordinates the lattice density is $s$. If $\ell$ is a
linear functional strictly positive on $\widetilde K\setminus\{0\}$, then
\begin{equation}\label{jet:count}
 \begin{aligned}
 &\#\{\gamma'\in\widetilde\Gamma:\ell(\gamma')<q\}\\
 &\qquad=s\,\operatorname{Leb}\{(D_0,x)\in\widetilde K:\ell(D_0,x)<1\}\,q^n+o(q^n),
 \end{aligned}
\end{equation}
with Lebesgue measure in the scaled coordinates. No saturation
assumption on $\Gamma$ is required.
\item
The vector
\begin{equation}\label{jet:canonical-vector}
 c_0=(r,\mathbf1)\qquad\text{in the original }(j,\gamma)
 \text{ coordinates}
\end{equation}
belongs to $\Lambda$. Its scaled coordinates are $(N,\mathbf1)$.
\end{enumerate}
\end{proposition}

\begin{proof}
Leading monomials multiply, and the associated leading coefficient is a
scalar. This proves additivity and the one-dimensional-leaf assertion.
In particular, each leading index in a fixed degree accounts for one
dimension of $S_j$.

Choose an integer $M>0$ divisible by all stabilizer orders such that
$ML$ descends to an ample line bundle on the coarse space. In a fixed
residue class $j_0$ modulo $M$, prescribe a finite compatible Taylor jet
of a section of $j_0L$. The kernel of the map to that finite jet space
is coherent. Serre vanishing after tensoring by sufficiently large
powers of $ML$, or after exact tame pushdown to the coarse space,
makes the map on global sections surjective. Thus the prescribed jet is
realized in every sufficiently large degree $j_0+kM$. Taking all the
finitely many admissible residue classes proves the assertion about
progressions. In particular, constants occur in all sufficiently large
degrees divisible by $m$, and each monomial $z^\gamma$ occurs as a
leading term in sufficiently large compatible degrees.

Differences of two such constant indices whose degrees differ by $m$
give $(m,0)$ in the generated group. For each $i$, a realized linear
jet gives an index $(j_i,e_i)$ with $j_i\equiv b_i\pmod m$.
The vectors $(m,0),(j_1,e_1),\ldots,(j_N,e_N)$ generate the full
congruence lattice in Equation~\eqref{jet:lattice}. All indices satisfy
that congruence, proving equality of groups. These same indices show
that the cone contains the positive constant ray, has full dimension,
and projects onto the exponent orthant.

We next prove the uniform estimate. Fix torus eigen-generators of $S$
and the resulting affine embedding of $C$. Their $\zeta_\nu$-degrees
are bounded below by a fixed positive multiple of $r_\nu$ for all
sufficiently large $\nu$, because their $\xi$-weights are positive.
A section of degree $j$ has a homogeneous polynomial lift in these
generators, of ordinary polynomial degree at most $Cj/r$. At a smooth
point of the punctured cone, a homogeneous function has the same local
vanishing order as its restriction to a transverse slice: in orbit and
slice coordinates it is a unit power of the orbit coordinate times its
slice expression.

Take a general complete-intersection curve through the smooth point,
cut out by affine hyperplanes in the fixed embedding. It is smooth
there and is not contained in the zero locus of the nonzero function.
The local intersection multiplicity with the function is at least its
local vanishing order. B\'ezout's inequality bounds this multiplicity
by the polynomial degree times the degree of the fixed embedded cone.
This proves Equation~\eqref{jet:bezout-bound} for $\Gamma$, and then for
$K$ by homogeneity and closure. In particular, bounded degree cuts are
compact.

The positive constant ray makes every nonempty fiber over an exponent
$x$ an upper ray in degree. Every exponent fiber is nonempty, since
nonnegative combinations of the indices $(j_i,e_i)$ reach every
$x\ge0$. Closedness and the degree bound make the lowest degree in each
fiber attainable. Hence $K$ is the epigraph of a finite convex
positively homogeneous function in the scaled coordinates.
Equation~\eqref{jet:bezout-bound} becomes
Equation~\eqref{jet:f-coercive}. A full-dimensional closed convex cone
is the closure of its interior.

The original lattice has covolume $m$. The coordinate change
$(j,\gamma)\mapsto(Nj/r,\gamma)$ multiplies covolumes by $N/r$;
the resulting covolume is $Nm/r=1/s$.
For clarity, we give the nonsaturated semigroup counting argument.
This is the conductor-and-exhaustion method for value semigroups
\cite[Theorems~1.4 and~1.6]{KavehKhovanskii2012}, also used for
graded linear series in \cite[Proposition~2.1]{LazarsfeldMustata2009}.
Its asymptotic conclusion does not assert exact saturation.
Let $\Gamma_0$ be a finitely generated subsemigroup with full generated
group $\Lambda$, and let $K_0$ be its cone. There is a vector
$c\in\Gamma_0$ such that
\begin{equation}\label{jet:conductor}
 c+(K_0\cap\Lambda)\subset\Gamma_0.
\end{equation}
Indeed every lattice point of $K_0$ lies in a cone generated by a
linearly independent subset of the finite generating set. Removing the
integer parts of its nonnegative coefficients writes it as a
semigroup element plus a lattice point in one of finitely many bounded
parallelepipeds. There are only finitely many possible remainders $v$.
Since the group is $\Lambda$, write each such remainder as
$v=v_+-v_-$ with $v_\pm\in\Gamma_0$. The sum of all the $v_-$ is a
single $c$ for which Equation~\eqref{jet:conductor} holds. Translating a
bounded linear cutoff changes its scale by a constant, so ordinary
lattice counting between $c+(K_0\cap\Lambda)$ and $K_0\cap\Lambda$
gives its leading lattice volume. After the coordinate change this
is Euclidean volume with density $s$.

Now exhaust $\Int K$ by cones of finitely many indices of $\Gamma$,
always retaining a fixed finite set whose group is $\Lambda$.
Every interior point belongs to the interior of one of these finite
cones: choose a small simplex surrounding it in $\Int K$ and
approximate the simplex vertices by elements of
$\operatorname{cone}(\Gamma)$. The resulting finite cones give the lower
leading count. Counting all lattice points of $K$ gives the upper
count. Compact convex cutoff bodies have boundary of Lebesgue measure
zero, and their volumes are exhausted from the interior. This proves
Equation~\eqref{jet:count}.

Finally, on the uniformizer the canonical character is the negative of
the sum of the coordinate characters, in the convention for equivariant
line bundles used here. The identity $-K_B=rL$ therefore gives
$r\equiv\sum_i b_i\pmod m$, proving
Equation~\eqref{jet:canonical-vector}.
\end{proof}

\subsection{Discrepancy and moment tests}

\begin{lemma}[The degree and Taylor-order test]\label{jet:test}
For $\delta\in\QQ_{\ge0}$, the rule
\begin{equation}\label{jet:test-rule}
 v_\delta\Bigl(\sum_j s_j\Bigr)
   =\min_{s_j\ne0}\left\{\frac{Nj}{r}
                           +\delta\ord_z(s_j)\right\}
\end{equation}
defines a valuation centered at the vertex. It is divisorial up to
positive rescaling, and
\begin{align}
 A_C(v_\delta)&=N+N\delta,\label{jet:test-discrepancy}\\
 A_C(v_\delta)^n\vol(v_\delta)
   &=N^n(1+\delta)^n s
        \int_{\RR_{\ge0}^N}e^{-f(x)-\delta|x|}\,dx.
                                      \label{jet:test-volume}
\end{align}
\end{lemma}

\begin{proof}
Use the local line-bundle extraction over the chosen uniformizer.
Write $t$ for its line coordinate, of grading weight one. Its finite
quotient presentation has the action
\[
 (t,z_1,\ldots,z_N)\longmapsto
 (\varepsilon^{-1}t,\varepsilon^{b_1}z_1,\ldots,
                                      \varepsilon^{b_N}z_N).
\]
A homogeneous function of degree $j$ pulls back to $t^jF_j(z)$.
The monomial valuation with costs
\[
 v(t)=N/r=1/(sm),\qquad v(z_i)=\delta
\]
therefore restricts to Equation~\eqref{jet:test-rule}. Distinct original
degrees have different powers of $t$, so even when scalar values agree
their lowest monomials cannot cancel. Rational positive costs give a
divisorial monomial valuation up to rescaling. At $\delta=0$ the rule is
the rescaled ordinary degree valuation. The positive degree cost and
the positivity of all algebra-generator degrees imply that its center
is the vertex.

Let $\Omega_C$ be the homogeneous canonical generator of weight $r$.
On this chart it has the form
\[
 t^{r-1}u(z)\,dt\wedge dz_1\wedge\cdots\wedge dz_N,
 \qquad u(0)\ne0.
\]
The exponent follows from homogeneity, and the coefficient is a unit
because $\Omega_C$ is nonvanishing on the punctured cone. Computing
discrepancy from this pulled-back form, including the finite-chart
ramification in the restricted value scaling, gives
\[
 A_C(v_\delta)=\frac Nr+N\delta+(r-1)\frac Nr
             =N+N\delta.
\]
Equivalently one can take a weighted monomial extraction on the smooth
chart and apply the canonical divisor formula there; the same
calculation descends to the quotient. For $\delta=0$ it is the canonical
degree formula.

In each $S_j$, a leading-term basis is compatible with total Taylor
order. Thus the colength below a cutoff $q$ is the number of indices
of $\Gamma$ satisfying $D_0+\delta|x|<q$. These finite quotients are
supported at the vertex, so computing them in $S$ or in its local ring
gives the same length. Proposition~\ref{jet:setup} applies to this
strictly positive cutoff functional. Since $K$ has dimension $n$,
homogeneity and integration by layers give
\[
 n!\operatorname{Leb}\{(D_0,x)\in\widetilde K:D_0+\delta|x|<1\}
 =\int_{\widetilde K}e^{-D_0-\delta|x|}\,dD_0\,dx
 =\int_{\RR_{\ge0}^N}e^{-f(x)-\delta|x|}\,dx.
\]
Combining this identity with the lattice density and
Equation~\eqref{jet:test-discrepancy} proves
Equation~\eqref{jet:test-volume}.
\end{proof}

Let $X_1,\ldots,X_N$ be independent exponential random variables of
mean one. For the current grading put
\begin{equation}\label{jet:moment-notation}
 B_0(x)=f(x)-|x|,\qquad
 H=\mathbb E e^{-B_0(X)}=\int_{\RR_{\ge0}^N}e^{-f(x)}\,dx.
\end{equation}

\begin{lemma}[The two moment inequalities]\label{jet:moments}
Along the fixed-cone sequence of gradings, there are numbers
$\eta_\nu\ge0$ tending to zero such that
\begin{equation}\label{jet:moment-inequalities}
 H\ge\frac2s,\qquad
 \mathbb E[(1+B_0(X))e^{-B_0(X)}]\ge-\eta_\nu H.
\end{equation}
For an exact minimizing rational grading one may take $\eta_\nu=0$.
\end{lemma}

\begin{proof}
Write $F(\delta)$ for the right side of
Equation~\eqref{jet:test-volume}. The volume hypothesis and the test at
zero give
\[
 2N^n\le\nvol(o,C)\le F(0)=N^n sH,
\]
which is the first inequality.

Consider the probability measure $d\mu_f=H^{-1}e^{-f(x)}dx$.
Positive homogeneity gives, for every $t>0$,
\[
 \int e^{-tf(x)}dx=t^{-N}H.
\]
Differentiating, justified by Equation~\eqref{jet:f-coercive}, yields
\begin{equation}\label{jet:homogeneous-moments}
 \mathbb E_{\mu_f}f=N,\qquad
 \mathbb E_{\mu_f}f^2=N(N+1).
\end{equation}
In particular,
\begin{align}
 \frac{F'(0)}{F(0)}
 &=n-\mathbb E_{\mu_f}|x|\notag\\
 &=1+\mathbb E_{\mu_f}(f-|x|)
 =\frac{\mathbb E[(1+B_0)e^{-B_0}]}{H}.
                                           \label{jet:first-derivative}
\end{align}

We check the uniformity needed for near-minimality.
Equation~\eqref{jet:f-coercive} and
Equation~\eqref{jet:homogeneous-moments} give uniform first and second
moments of $|x|$ under $\mu_f$. Indeed they are at most
$c_C^{-1}N$ and $c_C^{-2}N(N+1)$, respectively. Since
\[
 \frac{F(\delta)}{F(0)}
   =(1+\delta)^n\mathbb E_{\mu_f}e^{-\delta|x|},
\]
twice differentiating shows
\begin{equation}\label{jet:second-derivative}
 \sup_{0\le\delta\le1}\frac{|F''(\delta)|}{F(0)}\le C'_C
\end{equation}
for a fixed finite constant. The three terms are bounded by constant
multiples of $1$, $\mathbb E_{\mu_f}|x|$, and
$\mathbb E_{\mu_f}|x|^2$, because $e^{-\delta|x|}\le1$.

The test at zero is a rescaling of the current degree valuation, so
Lemma~\ref{cone:approximation} gives
$F_\nu(0)\to\nvol(o,C)$. If $F'_\nu(0)/F_\nu(0)$ were bounded above by
a fixed negative number along a subsequence, Taylor's formula and
Equation~\eqref{jet:second-derivative} would give a fixed small positive
rational $\delta$ and a fixed $\rho>0$ such that
$F_\nu(\delta)\le(1-\rho)F_\nu(0)$. This is incompatible with
$F_\nu(\delta)\ge\nvol(o,C)>0$ and the preceding convergence.
Thus $F'_\nu(0)/F_\nu(0)\ge-\eta_\nu$ for some $\eta_\nu\to0$.
Equation~\eqref{jet:first-derivative} proves the second inequality. At
an exactly minimizing degree valuation, every positive rational test
has volume at least $F(0)$, so its right derivative is nonnegative.
\end{proof}

\subsection{A uniform exponential expectation inequality}

We need to bound $H$ when $B_0$ is bounded below by a linear
function whose coefficients sum to zero. The second inequality in
Lemma~\ref{jet:moments} controls the expectation of $(1+B_0)e^{-B_0}$. Combining it with $H$
leads, for $u>0$, to $(1+(1+B_0)/u)e^{-B_0}$, whose pointwise upper envelope
is the truncated exponential below.

Define the continuously differentiable function
\begin{equation}\label{prob:g}
 g(y)=\begin{cases}
 1,&y\le0,\\
 (1+y)e^{-y},&y>0,
 \end{cases}
\end{equation}
Here $0\le g\le1$.

For any real $v$ and any $u>0$, elementary differentiation gives
\begin{equation}\label{jet:envelope}
 \sup_{b\ge v}\left(1+\frac{1+b}{u}\right)e^{-b}
       =\frac{e^u}{u}g(u+v).
\end{equation}
Indeed the derivative of the expression on the left before
maximization is $-(u+b)e^{-b}/u$, so its maximum is attained at
$b=-u$ when $v\le-u$, and at $b=v$ otherwise.

\begin{theorem}[The exponential bound]\label{prob:bound}
For $u=4/3$ one has
\begin{equation}\label{prob:universal-bound}
 \sup_{N\ge1}\ \sup_{\substack{a\in\RR^N\\\sum_i a_i=0}}
 \frac{e^u}{u}\mathbb E g(u+a\cdot X)
 =C_*:=\frac34e^{4/3}-\frac e3<2.
\end{equation}
The supremum over dimensions is approached by the vectors
$(2,-2/(N-1),\ldots,-2/(N-1))$ as $N\to\infty$.
\end{theorem}

The theorem gives a bound strictly below the threshold $2$ in
Lemma~\ref{jet:moments}. We prove it before applying it to the
canonical point. For the proof, put $h=1-g$; then
\begin{equation}\label{prob:kernel}
 h'(y)=ye^{-y}\mathbf1_{\{y>0\}},\qquad
 k(y):=h''(y)=(1-y)e^{-y}\mathbf1_{\{y>0\}}
\end{equation}
holds almost everywhere. In particular $h'$ is globally Lipschitz and
$k$ is bounded.

\begin{lemma}[Log-concavity and a barycenter half-space bound]
\label{prob:logconcavity}
A nonzero mean-zero linear combination of independent mean-one
exponentials has a positive continuous log-concave density on the whole
real line. Adding further independent scaled exponentials, with either
sign of coefficient, preserves these properties. If $Z$ has such a
density, then
\begin{equation}\label{prob:grunbaum}
 \frac1e\le\mathbb P(Z\le\mathbb EZ)\le1-\frac1e.
\end{equation}
\end{lemma}

\begin{proof}
The preservation of log-concavity under convolution is a special case
of the classical theory of Pr\'ekopa
\cite[Theorems~7--8]{Prekopa1973}; we give the elementary
exponential-specific argument needed here. A nonzero coefficient
vector with sum zero has a positive and a negative coefficient. The
difference $aE-bE'$ for $a,b>0$ has density
\[
 p(x)=\frac1{a+b}
 \begin{cases}e^{x/b},&x\le0,\\e^{-x/a},&x\ge0.
 \end{cases}
\]
It is positive, continuous, and log-concave.

Let $r$ be a positive log-concave function for which the indicated
one-sided integrals are finite. Set
$R(x)=\int_{-\infty}^x r(y)dy$. For $t\ge0$, concavity of $\log r$
implies that $r(x-t)/r(x)$ is nondecreasing in $x$. Thus
\[
 \frac{R(x)}{r(x)}=\int_0^\infty\frac{r(x-t)}{r(x)}dt
\]
is nondecreasing, and $(\log R)'=r/R$ is nonincreasing. Consequently
$R$ is log-concave. The same argument with $r(x+t)/r(x)$ shows that
$T(x)=\int_x^\infty r(y)dy$ is log-concave: $T/r$ is nonincreasing,
so $(\log T)'=-r/T$ is nonincreasing. These statements can equivalently
be read almost everywhere; the positive functions involved are locally
absolutely continuous.

If $p$ is a positive log-concave probability density, adding $cE$ or
$-cE$, $c>0$, produces respectively the densities
\[
 \frac{e^{-x/c}}c\int_{-\infty}^x e^{y/c}p(y)dy,
 \qquad
 \frac{e^{x/c}}c\int_x^\infty e^{-y/c}p(y)dy.
\]
The integrands are exponential tilts of $p$, hence log-concave, and
the appropriate one-sided integrals are finite. The preceding argument
proves log-concavity of the new densities. Positivity and continuity
are immediate from the displayed formulas. Induction starting with
$aE-bE'$ proves the assertions about finite exponential sums, and
translation does not affect them.

For completeness, Equation~\eqref{prob:grunbaum} is the
one-dimensional log-concave form of Gr\"unbaum's centroid inequality
\cite{Grunbaum1960}; see \cite[Lemma~5.4]{LovaszVempala2007} and
\cite[Section~3]{MNRY18}. It follows directly here from the same integral
argument. If $G$ is the distribution function of $Z$, both $G$ and
$1-G$ are log-concave. Since $G(Z)$ is uniform on $(0,1)$, Jensen's
inequality gives
\[
 \log G(\mathbb EZ)\ge\mathbb E\log G(Z)
 =\int_0^1\log t\,dt=-1.
\]
Applying the same reasoning to $1-G$ gives
$\log(1-G(\mathbb EZ))\ge-1$. All the means in the present application
are finite, and both logarithmic expectations just used equal $-1$.
These two inequalities give Equation~\eqref{prob:grunbaum}.
\end{proof}

\begin{lemma}[Strict single crossing]\label{prob:single-crossing}
Let $Z$ have a positive log-concave probability density on $\RR$, and
let $E$ be an independent mean-one exponential. If
\[
 \mathbb E k(Z+c_*E)=0,
\]
then $\mathbb E k(Z+cE)>0$ for every $c<c_*$ and
$\mathbb E k(Z+cE)<0$ for every $c>c_*$.
\end{lemma}

\begin{proof}
Write $p$ for the density of $Z$ and first consider
\[
 \psi(t)=\mathbb E k(Z+t)
        =\int_0^\infty(1-y)e^{-y}p(y-t)dy.
\]
This is continuous, by continuity of translation in $L^1$ and
boundedness of $k$. A positive log-concave density has a continuous
representative whose logarithm is finite and concave. For $t_1<t_2$,
the ratio
\[
 r(y)=p(y-t_2)/p(y-t_1)
\]
is nondecreasing in $y$. If $\psi(t_1)=0$, subtraction of the constant
$r(1)$ gives
\[
 \psi(t_2)
 =\int_0^\infty k(y)p(y-t_1)(r(y)-r(1))dy\le0:
\]
the integrand is nonpositive on both $(0,1)$ and $(1,\infty)$.
Equality would force $r$ to be constant on $(0,\infty)$, by positivity
and continuity. Thus $\log p$ would have constant increments across
the fixed interval length $t_2-t_1$ on a half-line. Concavity forces
it to be affine there: its almost-everywhere derivative is
nonincreasing, and equality of this derivative across a fixed shift
makes it constant on that half-line. Integrability then gives
$p(y-t_1)=A e^{-\lambda y}$ for $y>0$, with $A,\lambda>0$. But
\[
 \psi(t_1)=A\int_0^\infty(1-y)e^{-(1+\lambda)y}dy
         =\frac{A\lambda}{(1+\lambda)^2}>0,
\]
a contradiction. Hence the inequality is strict. Reversing the
likelihood-ratio comparison proves strict positivity to the left of a
zero. Continuity shows that $\psi$ either has one strict sign
everywhere or has exactly one zero, with positive sign to its left
and negative sign to its right.

Now put $F(c)=\mathbb E\psi(cE)$. For $c\ne0$, the density of $cE$ is
\[
 q_c(t)=|c|^{-1}e^{-t/c}\mathbf1_{\{tc>0\}}.
\]
If $c_1<c_2$ have the same sign, then on their common support
\[
 \frac{q_{c_2}(t)}{q_{c_1}(t)}
 =\frac{|c_1|}{|c_2|}
             \exp\bigl(t(1/c_1-1/c_2)\bigr)
\]
is strictly increasing in $t$. This holds on the negative half-line
as well as on the positive half-line. If $F(c_*)=0$ and $c_*\ne0$,
then $\psi$ must have its strict crossing inside that half-line.
Subtracting the likelihood ratio at this crossing proves strict
positivity of $F(c)$ for $c<c_*$ and strict negativity for $c>c_*$
within the same parameter half-line.

The comparisons across zero are also strict. If $c_*>0$, the crossing
of $\psi$ lies in $(0,\infty)$, so $\psi>0$ on $(-\infty,0]$ and
$F(c)>0$ for $c\le0$. If $c_*<0$, the crossing lies in
$(-\infty,0)$, giving $F(c)<0$ for $c\ge0$. Finally $F(0)=\psi(0)$;
a zero there gives the two required signs directly from those of
$\psi$. This proves the assertion for all real parameters.
\end{proof}

\begin{lemma}[Finite critical points]\label{prob:critical-points}
Fix $u>0$ and $N\ge2$. Every nonzero finite local maximum of
\[
 a\longmapsto\mathbb E g\left(u+\sum_{i=1}^N a_iX_i\right),
 \qquad \sum_i a_i=0,
\]
has, after permutation, the form
\begin{equation}\label{prob:exceptional-vector}
 a=(t,-t/(N-1),\ldots,-t/(N-1))\qquad(t>0).
\end{equation}
\end{lemma}

\begin{proof}
Equivalently consider a local minimum of $\mathbb Eh(Y)$, where
$Y=u+\sum_i a_iX_i$. At a nonzero coefficient vector,
Lemma~\ref{prob:logconcavity} supplies a continuous positive
log-concave density for $Y$. First and second differentiation under the
expectation are justified by boundedness of $h'$, its global Lipschitz
property, the finite second moments of the variables $X_i$, and the absence of an atom at $Y=0$.
Stationarity in the direction $e_i-e_j$ gives
\begin{equation}\label{prob:stationarity}
 0=\mathbb E[(X_i-X_j)h'(Y)]
   =(a_i-a_j)\mathbb E[X_iX_jk(Y)].
\end{equation}
To see the second equality, condition on all other variables and on
$S=X_i+X_j$. Given $S=s$, $X_i$ is uniform on $[0,s]$ and
$X_j=s-X_i$. Integration by parts with $x(s-x)$, whose endpoint
values vanish and whose derivative is $s-2x$, gives exactly
Equation~\eqref{prob:stationarity}.

For distinct indices, weighting by $X_iX_j$ changes those two
independent exponential laws into independent Gamma$(2,1)$ laws. Since
$\mathbb E X_iX_j=1$, no normalization factor appears. Thus
\begin{equation}\label{prob:size-bias}
 \mathbb E[X_iX_jk(Y)]
   =\mathbb E k(Y+a_iE'+a_jE''),
\end{equation}
where the added exponentials are independent of all variables in $Y$.
If three distinct coefficient values occurred, fix an index $i$ at
one value and choose indices at the other two. Equations
\eqref{prob:stationarity} and \eqref{prob:size-bias} would give two
distinct zeros of
\[
 c\longmapsto\mathbb E k(Y+a_iE'+cE'').
\]
Its base variable has a positive full-support log-concave density by
Lemma~\ref{prob:logconcavity}, contradicting
Lemma~\ref{prob:single-crossing}. Therefore there are at most two
coefficient values.

Suppose the two values are $b<c$ and the higher value occurs at
indices $i\ne j$. A high-low pair gives a zero at $b$ of
$d\mapsto\mathbb E k(Y+cE'+dE'')$. Strict single crossing makes its
value at $c$ negative. By Equation~\eqref{prob:size-bias}, this says
$\mathbb E[X_iX_jk(Y)]<0$. As $a_i=a_j$, $Y$ depends on this pair only
through $S$, and
\[
 \mathbb E[(X_i-X_j)^2\mid S]=S^2/3,
 \qquad \mathbb E[X_iX_j\mid S]=S^2/6.
\]
Consequently the Hessian in the permissible splitting direction is
\[
 \mathbb E[(X_i-X_j)^2k(Y)]
       =2\mathbb E[X_iX_jk(Y)]<0,
\]
contradicting a minimum. The higher value therefore occurs once.
The sum-zero constraint forces it to be positive and gives
Equation~\eqref{prob:exceptional-vector}.
\end{proof}

The preceding lemmas leave only one possible shape for a finite
nonzero optimizer. To obtain a bound uniform in the dimension, it
remains to control escape to infinity and then solve the resulting
one-variable maximization.



\begin{proof}[Proof of Theorem~\ref{prob:bound}]
Put $\Phi_N(a)=(e^u/u)\mathbb E g(u+a\cdot X)$ and let
$\mathcal P_N$ be its supremum under the sum-zero constraint. First
control escape to infinity with $N$ fixed. From an unbounded sequence
$a_k$, pass to a subsequence with
$a_k/\|a_k\|\to b\ne0$. The vector $b$ has sum zero, so $b\cdot X$
has mean zero and a continuous log-concave density. Almost surely
away from its zero set,
\[
 g(u+a_k\cdot X)\longrightarrow\mathbf1_{\{b\cdot X<0\}}.
\]
Bounded convergence and Equation~\eqref{prob:grunbaum} imply that every
such subsequential limit is at most
\begin{equation}\label{prob:boundary-value}
 B_*:=\frac{e^u}{u}(1-1/e).
\end{equation}
Thus, if $\mathcal P_N>B_*$, a maximizing sequence is bounded and a
finite maximizer exists by continuity. The zero vector has value
$1+1/u=7/4$. Every other finite maximizer has the form in
Lemma~\ref{prob:critical-points}. If $\mathcal P_N\le B_*$, the boundary
estimate already suffices; we make no assertion of attainment in
that case.

Appending a zero coefficient shows that $\mathcal P_N$ is
nondecreasing in $N$, and it is bounded by $e^u/u$. Consider its limit.
If that limit is at most $\max(B_*,7/4)$, those two values suffice.
Otherwise sufficiently large dimensions have finite nonzero
maximizers. After permutation write them as in
Equation~\eqref{prob:exceptional-vector}, with exceptional coefficient
$t_N>0$. On a common probability space their arguments are
\[
 Y_N=u+t_N\left(X_1-\frac1{N-1}\sum_{i=2}^NX_i\right).
\]
Pass to a subsequence on which $t_N$ converges in $[0,\infty]$.
If $t_N\to t<\infty$, the law of large numbers and bounded convergence
give
\[
 \Phi_N(a_N)\longrightarrow
 \Phi_\infty(t):=\frac{e^u}{u}\mathbb E g(u+t(X_1-1)).
\]
If $t_N\to\infty$, then $Y_N/t_N\to X_1-1$ almost surely. The limit
has no atom at zero; on its positive and negative sets, respectively,
$Y_N$ tends to $+\infty$ and $-\infty$. Hence bounded convergence gives
\[
 \Phi_N(a_N)\longrightarrow
       \frac{e^u}{u}\mathbb P(X_1<1)=B_*.
\]
This argument uses no comparison between the growth of $t_N$ and
$\sqrt N$.

It remains to maximize the scalar expression. For $0\le t\le u$,
the argument $u+t(X_1-1)$ is nonnegative, and direct integration gives
\begin{equation}\label{prob:scalar-formula}
 \Phi_\infty(t)=
 \begin{cases}
 \displaystyle\frac{e^t}{u}
       \left(\frac{1+u-t}{1+t}+\frac{t}{(1+t)^2}\right),
                                      &0\le t\le u,\\[7pt]
 \displaystyle\frac{e^u}{u}
       \left(1-e^{-1+u/t}\frac{t^2}{(1+t)^2}\right),
                                      &t\ge u.
 \end{cases}
\end{equation}
For the second branch, split the exponential integral at
$X_1=1-u/t$; above that point translate the integration variable so
that the argument of $g$ is $t$ times the new variable. The integral
there is
$e^{-1+u/t}((1+t)^{-1}+t(1+t)^{-2})$, which yields the displayed
formula.

The derivative of the first branch has the sign of
$t(-t^2+ut+u-1)$. For $u=4/3$, the second factor is positive throughout
$[0,u]$, as its endpoint values are positive and it is concave. On the
second branch, the logarithmic derivative of the positive quantity
being subtracted is
\[
 -\frac u{t^2}+\frac2t-\frac2{1+t}
       =\frac{(2-u)t-u}{t^2(1+t)}.
\]
Thus this branch is maximal at $t=u/(2-u)=2$, and its value is
\[
 \Phi_\infty(2)=\frac34e^{4/3}-\frac e3=C_*.
\]
In particular $C_*>\Phi_\infty(0)=7/4$. Moreover
\[
 C_*-B_*=\frac34e^{1/3}-\frac e3>0,
\]
because $e^{1/3}>1+1/3$ and $e<3$. The preceding compactness and
limit argument proves the upper bound $\sup_N\mathcal P_N\le C_*$.
Taking $t_N=2$ and applying the law of large numbers proves the reverse
inequality and the asserted approach to the supremum.

Finally the strict gap has an elementary rational certificate. The
exponential series gives $e>1359/500$ and $e^{4/3}<1897/500$.
For the latter bound one may retain terms through degree eight and
bound the remaining tail by its first term divided by
$1-(4/3)/10$. Therefore
\[
 C_*<\frac34\frac{1897}{500}-\frac13\frac{1359}{500}
     =\frac{3879}{2000}<2.
\]
\end{proof}

\subsection{Interiority of the canonical vector}

\begin{theorem}[Canonical interior point]\label{jet:interior}
For every sufficiently large grading in the subsequence under
consideration,
\begin{equation}\label{jet:interior-conclusion}
 c_0=(r,\mathbf1)\in\Int K
                 \qquad\text{in the original degree coordinates}.
\end{equation}
Equivalently, $f(\mathbf1)<N$. For an exact minimizing rational grading
with $m\ge r/N$, the conclusion holds without passage to a limit.
\end{theorem}

\begin{proof}
Suppose that $f(\mathbf1)\ge N$. At the interior point $\mathbf1$ of
the exponent orthant, choose a subgradient $p$ of $f$. Positive
homogeneity gives $p\cdot\mathbf1=f(\mathbf1)$ and
$f(x)\ge p\cdot x$ on the orthant. Therefore
$B_0(x)\ge(p-\mathbf1)\cdot x$, and the coefficient sum of
$p-\mathbf1$ is nonnegative. Subtracting the same nonnegative number
from each coefficient gives a vector $a$ with sum zero and
\begin{equation}\label{jet:supporting-plane}
 B_0(x)\ge a\cdot x\qquad(x\in\RR_{\ge0}^N).
\end{equation}

Take $u=4/3$. Applying Equation~\eqref{jet:envelope} pointwise to
Equation~\eqref{jet:supporting-plane}, and using
Lemma~\ref{jet:moments}, gives
\begin{align*}
 (1-\eta_\nu/u)H
 &\le\mathbb E\left[\left(1+\frac{1+B_0(X)}u\right)
                                      e^{-B_0(X)}\right]\\
 &\le\frac{e^u}{u}\mathbb E g(u+a\cdot X)
 \le C_*<2.
\end{align*}
All expectations are finite by Equation~\eqref{jet:f-coercive}.
But $H\ge2/s\ge2$, so the displayed inequality is impossible once
$\eta_\nu$ is sufficiently small. Hence $f(\mathbf1)<N$ for all
sufficiently large gradings. At an exact minimizing rational grading,
$\eta_\nu=0$ gives the contradiction immediately.

The scaled coordinates of $c_0$ are $(N,\mathbf1)$, and
$\mathbf1$ is interior to the exponent orthant. The strict epigraph
inequality is therefore precisely
Equation~\eqref{jet:interior-conclusion}. Its membership in the lattice
was already proved in Proposition~\ref{jet:setup}.
\end{proof}
